\chapter{Programmdokumentation}

Das CZGame ist ein am Carl-Zeiss-Gymnasium Jena orientiertes, als Point- \& Click ausgelegtes Spiel. Ziel ist es sich gegen alle Lehrer aus den verschiedenen Fachschaften durchzusetzen. Dabei wird man im Verlauf des Spiels immer besser und bekommt immer stärkere Attribute. \\

\begin{figure}[h!]
    \centering
    \includegraphics[width=0.5\linewidth]{bilder/StartSzene.png}
    \caption{Die Startszene des CZGame}
    \label{fig:placeholder}
\end{figure}

Das Spiel startet in einer Szene, die dem Foyer der Schule ähnelt. Von hier aus kann man sich mithilfe der Pfeile fortbewegen. Man kann entweder nach links und rechts um in die verschiedenen Gänge der Fachschaften zu gelangen, oder man geht die Treppe in die nächste Etage hinauf. In den Gängen der Fachschaften kann man nun Minigames finden. Diese beinhalten 3 unterschiedlich schwere Level (leicht, mittel, schwer). Die Minigames sind an die jeweilige Fachschaft angelegt. Beispielsweise muss man in der Informatik Logische Schaltungen durchblicken, während der Spieler in der Chemie mit dem sortieren von verschiedenen Flüssigkeiten konfrontiert wird. Schafft man es die Minigames erfolgreich zu lösen, erhält man verschiedenste Items, die einem im Kampf gegen die Lehrer unterstützen. \\
In den Gängen kann man außerdem in die Räume gehen. Dafür genügt es, wenn man die Tür anklickt. In den Räumen findet man die verschieden Lehrer. Möchte man sich im Kampf gegen diese Beweisen muss man einmal auf den Lehrer klicken. \\
Der Kampf basiert auf verschiedenen Runden. Lehrer und Schüler versuchen abwechseln sich ihrem gegenüber zu beweisen. Durch die Items kann man stärker werden und die Lehrer leichter besiegen. Dabei muss der Spieler jedoch aufpassen, denn nicht jedes Item hilft bei jedem Lehrer gleich viel. 


\section{Engine}
\input{texte/04_programmdoku-Engine.tex}

\section{Grafikdesign}
\input{texte/04_programmdoku-Grafikdesign.tex}

\section{Kampfkonzept}
\input{texte/04_programmdoku-Kampfkonzept.tex}

\section{Minigames}
\input{texte/04_programmdoku-Minigames.tex}

\section{Item- und Charakterdesign}
\input{texte/04_programmdoku-Item_und_Characterdesign.tex}

